\subsection{Dimension}

\begin{exer}
	Let $X$ be a topological space. Let $\{X_i\}_i$ be a locally finite
	covering of $X$ by closed subsets $X_i$. Then $\dim X = \sup_i \dim
	X_i$.
\end{exer}
\begin{proof}
	Let $x$ be a point of $X$ of dimension $n$. Choosing a neighborhood $U$
	of $x$ that only meets finitely many of the $X_i$, let $Z_0 \subsetneq
	Z_1 \subsetneq \cdots \subsetneq Z_n$ be an ascending chain of
	irreducible subsets. If $X_1,\ldots,X_m$ are the $X_i$ that meet $Z_n$,
	it is clear that at least one of the $X_i$ contains $Z_n$, hence
	$\sup_i \dim X_i \geq n$ and the proposition holds.
\end{proof}

\begin{exer}
	Let $X$ be a scheme and $Z$ be a closed subset of $X$. Then for all $x
	\in X$, we have $\codim(\overline{\{x\}}, X) = \dim \Sh_{X,x}$ and
	$\codim(Z,X) = \min_{z \in Z} \dim \Sh_{X,z}$.
\end{exer}
\begin{proof}
	Let $\overline{\{x\}} \subseteq Z_0 \subseteq \cdots \subseteq Z_n$ be
	an ascending chain of irreducible closed subsets of $X$. This induces a
	chain of specializations $z_n \rightsquigarrow \cdots \rightsquigarrow
	z_0 \rightsquigarrow x$, so all of the $z_i$ lie in the image of
	$\Spec \Sh_{X,x}$ and $\dim \Sh_{X,x} \geq \codim(\overline{\{x\}},X)$.
	The other direction is obvious.\\

	The second statement follows from the same argument as above.
\end{proof}

\begin{exer}
	We have the following:
	\begin{enumerate}[label=(\alph*)]
		\item Let $Z$ be a closed subset of a topological space $X$.
			Then $\codim(Z,X) = 0$ iff $Z$ contains an irreducible
			component of $X$. This does not mean that $\dim Z$ is
			equal to $\dim X$.
		\item Let $X = \Spec \Dvr_K[T]$, where $\Dvr_K$ is a discrete
			valuation ring with uniformizing paramater $t \neq 0$.
			Let $f = tT - 1$. Then $(f)$ is a maximal ideal. Let
			$x \in X$ be the corresponding point. Then $X$ is
			irreducible, $\dim \Sh_{X,x} = 1$, and $\codim(\{x\},X)
			+ \dim \{x\} < \dim X$.
	\end{enumerate}
\end{exer}
\begin{proof}
	Part (a) is obvious. For part (b), let $a_0 \in \Dvr_K$ be contained in
	$(t)$. Since $a_0 = ut^n$ for some unit $u$, we have
	$1 = (1 + \cdots + tT^{n-1})(1-tT) + u^{-1}a_0$ and so $(f, a_0)$ is the
	unit ideal. The case where $a_0$ is replaced with a general polynomial
	follows easily, so $(f)$ is a maximal ideal. That $X$ is irreducible is
	a consequence of the fact that $\Dvr_K[T]$ is entire. To see that
	$\Sh_{X,x}$ is of dimension 1, just notice that it is a discrete
	valuation ring. Finally, we have that $\codim(\{x\},X) + \dim(\{x\}) =
	1$, but $\dim(\Dvr_K[T]) = 2$.
\end{proof}

\begin{exer}
	Let $A$ be a ring and $\mathfrak p_1,\ldots,\mathfrak p_r$ be prime
	ideals of $A$. Let $I$ be an ideal of $A$ contained in none of the
	$\mathfrak p_i$. Then $I$ is not contained in $\mathfrak p_1 \cup \cdots
	\cup \mathfrak p_r$.
\end{exer}
\begin{proof}
	Suppose $r = 2$. Then if $x \in I \setminus \mathfrak p_1$, we should
	assume that $x \in \mathfrak p_1$. Let $y$ be in $I \setminus \mathfrak
	p_2$. If $x + y \in \mathfrak p_1$, $y \not\in \mathfrak p_1$ and $y$
	is not contained in either ideal. Otherwise, $x \in \mathfrak p_1$,
	contradiction.\\

	Now, assume that the statement holds for $r-1$ ideals. Letting $x_i$ be
	an element of $I \setminus (\bigcup_{j \neq i} \mathfrak p_j)$, we see
	that $x_1 + x_2 \cdots x_r$ is not contained in any of the ideals.
\end{proof}

\begin{exer}
	Let $A$ be a graded ring and let $\mathfrak p_1,\ldots,\mathfrak p_r$ be
	homogeneous prime ideals of $A$, none of which contain $A_+$.
	Let $I$ be a homogeneous ideal generated by homogeneous elements of
	positive degree, contained in none of the $\mathfrak p_i$. Then there is
	a homogeneous element of $I$ not contained in $\mathfrak p_1 \cup \cdots
	\cup \mathfrak p_r$.
\end{exer}
\begin{proof}
	Suppose $r = 2$. If $x \in I \setminus \mathfrak p_1$ is homogeneous,
	suppose that $x \in \mathfrak p_2$. Let $y$ be in $I \setminus \mathfrak
	p_2$. Since both $x$ and $y$ are of non-zero degree, $x^{\deg y} +
	y^{\deg x}$ is a homogeneous element. It is easy to see that either
	$y$ or the sum above is the desired element.\\

	Assume now that the statement holds for $r-1$ ideals. Choosing a
	homogeneous element $x_i$ of $I \setminus (\bigcup_{j \neq i} \mathfrak
	p_j)$ for each $i$, we see that $x_1^{\deg x_2 + \cdots + \deg x_r} +
	(x_2 \cdots x_r)^{\deg x_1}$ is the desired element.
\end{proof}
